\section{Effective buffered caustics of the physical collision graph}
\label{sec:v63-effective-caustics}
We quantify the dependence on the collision count and on the incidence
cap in the geometric localization of
Section~\ref{sec:v62-clearance-caustics}. The reference set used here
is a specified algebraic overcover of the physical caustic. It carries
no measure. The source localized near it will remain a restriction of
the original orbit source. A factor two between the source cap and the
reference cap is needed when the cap varies.

\subsection{A finite-format reference graph}
Choose an integer $T_0$ larger than the uniform flight horizon. Use
compact rational angular charts with algebraic frames and introduce
$c_i>0$ by $c_i^2+p_i^2=1$. For each word, mark $j+1$, candidate disk
and chart, let $\widehat{\mathcal K}_{m,\eta}$ be the weak-first-hit
graph used in the counting argument of
Theorem~\ref{thm:v62-caustic-finite}. More precisely, retain the contact
and reflection equations, the chosen entry-root signs, all first-jet
equations, $\ell_0\le\ell_i\le T_0$, and $c_i\ge\eta$ for
$0\le i\le m$, and replace only the strict first-hit inequalities by
their weak versions. Retain also the compact selected-witness envelope
\[
 R\le |Z|\le R+s_0,\qquad L_d\ge1/5.
\]
This envelope contains the original one in module 134, since
$\ell_*>1/5$. Take the finite disjoint union over the indicated labels.
The full roof $F$, the clearance $Z$ and
$\mathcal J=\det D(F,Z)$ have exactly the marked-chart meanings in
\eqref{eq:v62-fiber-objects}. In particular $F$ contains the entire
backward and forward parts of the original $m$ flights.

Define
\begin{equation}\label{eq:v63-overcover}
 \widehat{\mathcal C}_{m,\eta}
 =\{(R,F(X)):X\in\widehat{\mathcal K}_{m,\eta},\quad
                      Z(X)^2=R^2,\quad\mathcal J(X)=0\}.
\end{equation}
The physical caustic $\mathcal C_{m,\eta}$ is contained in this set.
Equality is neither needed nor asserted. In contrast with the earlier
use of the weak graph solely for counting, its image now serves as a
larger geometric reference set. No trajectory or probability is put on
its additional points.

\begin{lemma}[Format and finite fibers of the reference graph]
\label{lem:v63-graph-format}
For each $m\ge2$ and $0<\eta\le1$ the graph and its critical image
are compact. At a fixed radius the image has at most $CA^m$ points.
For fixed $R,\eta$ and $0\le h\le1$, the compact set
\begin{equation}\label{eq:v63-horizontal-image}
 V_{m,\eta}(R,h)
 =\{t:(R',t)\in\widehat{\mathcal C}_{m,\eta},\ |R'-R|\le h\}
\end{equation}
is a union of at most $CA^m$ closed intervals and points.

There are fixed integers $a,d\ge2$ for which the following format
bounds hold. One word graph uses at most $a(m+1)$ real variables and
polynomial atoms of degree at most $d$. The union uses at most
$\exp(a(m+1))$ atoms after Boolean expansion, with the same variable
blocks. Coefficient bit lengths, after clearing rational denominators,
are at most $a(m+1)^2$. The only fixed algebraic coefficient needed is
$\sqrt3$, which may be replaced by one positive variable satisfying
$z^2=3$. The variables $R,\eta,h$ are not treated as coefficients.
\end{lemma}
\begin{proof}
This is the auxiliary physical graph of
Sections~\ref{sec:v61-inverse-incidence} and
\ref{sec:v62-clearance-caustics}, with first jets as in
Lemma~\ref{lem:v62-jet-complexity}. Intermediate contacts are retained,
not eliminated. Each new collision adds a fixed number of positions,
velocities, flight lengths and first-derivative variables. Introducing
a reciprocal for each nonzero implicit coefficient keeps the polynomial
degree bounded. Under a positive cap these coefficients do not vanish.
Their reciprocals and the unique first jets are bounded on each capped
compact parameter set. Thus the weak graph is closed and bounded.
These bounds may depend on $m,\eta$; no norm bound on the jets is
being claimed independent of them.

The next-center alphabet is fixed by $T_0$. In local coordinates its
coefficients lie in $\mathbb Q(\sqrt3)$ and have a fixed height bound.
Alternatively, absolute lifted centers have size $O(m+1)$, which is
still within the stated bit bound. Rational angular frames may be
chosen among the coordinate-axis frames; positive denominators are
retained as graph variables. Neither the moving section nor any
exact-label indicator is encoded here. There are exponentially many
words and only polynomially many marks. Their Boolean disjunction
reuses the same continuous variables: it does not introduce one
quantified block for each word. This proves the format statements.

For a fixed $R$, $dZ\ne0$ on the envelope: in the collision coordinates
$\partial_\varphi Z=L_d\ge1/5$. On the seam and rank locus,
$dF\wedge dZ=0$. The implicit simple-root equations define the same
smooth local jets even where a weak first-hit inequality becomes an
equality. On every smooth stratum of a path in that locus, $dZ=0$
and hence $dF=0$ along the path. Semialgebraic connected components
are piecewise smooth path connected. Continuity of $F$ at successive
stratum endpoints makes it constant on the entire component. This
includes zero-dimensional strata and junctions; no differentiability
across a change of selected word is used, since word copies are disjoint.

For each word, the sign-component estimate
\eqref{eq:v44-sign-bound} applies to $O(m+1)$ variables and atoms of
fixed degree, also after adding the two horizontal-strip inequalities
in \eqref{eq:v63-horizontal-image}. It gives $CD^m$ components.
Summing words and marks gives $CA^m$. For a fixed radius each critical
component contributes one value. In the horizontal strip its roof
projection is a compact connected subset of the real line, and thus
an interval or point. These observations prove both counting assertions.
\end{proof}

\subsection{An effective one-variable lower modulus}
We record the elementary scalar step used to keep the cap dependence
quantitative. Its algebraic input is effective quantifier elimination,
not a billiard mixing estimate.

\begin{lemma}[A positive monotone modulus of controlled format]
\label{lem:v63-effective-modulus}
Suppose $a_m:(0,1]\to(0,1]$ is positive and nondecreasing. Suppose its
graph has a first-order description over the reals with at most twelve
quantifier blocks, at most $a(m+1)$ variables in each block, and at
most a fixed number of free variables. Assume the degree, atom-count
and integer coefficient bounds of Lemma~\ref{lem:v63-graph-format}.
Then, for a fixed $C_0\ge1$ depending only on this format,
\begin{equation}\label{eq:v63-budget}
 E_m=\left\lceil2^{C_0(m+1)^{16}}\right\rceil,\qquad K_m=2^{E_m}
\end{equation}
satisfy
\begin{equation}\label{eq:v63-scalar-modulus}
 a_m(\delta)\ge K_m^{-1}\delta^{E_m},\qquad0<\delta\le1.
\end{equation}
The same $C_0$ can cover any fixed finite list of such descriptions.
\end{lemma}
\begin{proof}
The quantifier-elimination bound in \cite[Theorem 4]{BMN} gives a
quantifier-free description of the graph by integer polynomials whose
degrees $D$ and coefficient bit lengths $H$ satisfy
\[
 D,H\le2^{C_1(m+1)^{12}}.
\]
Indeed the degree bound is $d^{O(k_1)\cdots O(k_\omega)}$, and the
bit bound is the input bit length times that quantity, with
$\omega\le12$ and $k_i=O(m+1)$. The number of atoms can be large,
but does not enter these two bounds. Adjoining the fixed algebraic
coefficient can be merged into an existential block; increasing the
fixed block budget to twelve includes it.

Here is a direct passage from these bounds to a lower modulus; it
requires neither continuity of $a_m$ nor a quantitative Puiseux theorem.
At every graph point $(\delta,u)$ at least one nonzero defining
polynomial $P(\delta,u)$ vanishes. Otherwise the truth values of all
its signs are constant on a neighborhood, contradicting that the
formula defines the graph of a single-valued function. Remove the
largest power of $u$ dividing this polynomial and write
$P=P_0(\delta)+\sum_{j=1}^D P_j(\delta)u^j$, with $P_0\ne0$.
For
\[
 0<\delta\le\delta_0=[4(D+1)2^H]^{-1}
\]
the first nonzero integer coefficient of $P_0$ gives
$|P_0(\delta)|\ge\delta^D/2$. The sum of the absolute values of
all higher terms is at most $D(D+1)2^H u$ for $0<u\le1$.
Consequently every such positive root satisfies
\[
 u\ge\frac{\delta^D}{2D(D+1)2^H}.
\]
A polynomial independent of $u$ cannot vanish in this range, by the
same leading-coefficient estimate. These estimates hold for all the
polynomials at once, without multiplying their number.

For $\delta\ge\delta_0$, use monotonicity and the estimate at
$\delta_0/2$. It follows that the whole interval has a lower bound
$c\delta^D$, where $\log_2(c^{-1})\le C_2D(H+\log(2+D))$.
Increasing $C_0$ in \eqref{eq:v63-budget} dominates this last quantity
and $D$, and proves \eqref{eq:v63-scalar-modulus}. All constants are
controlled by the displayed finite format, not by an unrecorded
real coefficient of a fixed-cap distance function.
\end{proof}

\subsection{Buffered separation, uniform in the cap}
For $X\in\widehat{\mathcal K}_{m,\chi}$, $0<\chi\le1$, put
\begin{equation}\label{eq:v63-buffered-distance}
 d_{m,\chi}(X)=\min\{1,\operatorname{dist}((R,F(X)),
                                      \widehat{\mathcal C}_{m,\chi/2})\},
 \qquad f(X)=\bigl||Z(X)|-R\bigr|+|\mathcal J(X)|.
\end{equation}
The distance is in parameter--roof space. It is one if the reference
set is empty. All inequalities below apply in particular to the
original physical graph capped at $\chi$.

\begin{theorem}[Effective separation with a buffered incidence cap]
\label{thm:v63-buffered-separation}
After one fixed increase of $C_0$ in \eqref{eq:v63-budget},
\begin{equation}\label{eq:v63-buffered-separation}
 (\chi d_{m,\chi}(X))^{E_m}\le K_m f(X)
       \quad(m\ge2,\ 0<\chi\le1,\ X\in\widehat{\mathcal K}_{m,\chi}).
\end{equation}
Thus an original first-clearance point satisfies
\begin{equation}\label{eq:v63-rank-threshold}
 4K_m\varepsilon\le(\chi h)^{E_m},\quad d_{m,\chi}\ge h
 \quad\Longrightarrow\quad
 |\mathcal J|\ge\frac{(\chi h)^{E_m}}{2K_m}.
\end{equation}
\end{theorem}
\begin{proof}
For $0<\delta\le1$, let $a_m(\delta)$ be the infimum of the set
containing $1$ and every $\min(1,f(X))$ such that
\[
 \delta\le\chi,h\le1,\quad X\in\widehat{\mathcal K}_{m,\chi},
 \quad\operatorname{dist}((R,F(X)),\widehat{\mathcal C}_{m,\chi/2})\ge h.
\]
It is nondecreasing. It is strictly positive. To prove the latter,
suppose for a fixed $\delta>0$ that $f(X_l)\to0$ along eligible
points. Pass to a subsequence in one word/chart copy with
$\chi_l\to\chi_*$ and $X_l\to X_*$. The cap $\delta$ bounds the
simple-root jets, so such a subsequence exists. The limit is a critical
seam point with $c_i(X_*)\ge\chi_*$ for every $i$. For all sufficiently
large $l$ it belongs to the reference graph with cap $\chi_l/2$.
Its image approaches $(R_l,F(X_l))$, contradicting the distance lower
bound $\delta$. The buffer is used precisely at this step; continuity
of an unbuffered cap-dependent distance is not assumed.

The graph of this infimum has controlled format. Distance at least $h$
is the universal statement that every critical graph point $Y$ at cap
$\chi/2$ has squared image distance at least $h^2$. The empty-reference
case is included by that statement. The condition that $u$ is the
infimum is the conjunction of: $u\le1$ and $u$ is a lower bound for
all eligible points; and, for every $r>0$, either $1<u+r$ or there is
an eligible point with $\min(1,f)<u+r$. The first conjunct has prefix
$\forall X\,\exists Y$ after negating eligibility; the second has
prefix $\forall r\,\exists X\,\forall Y$. Caps, thresholds and the
absolute-value graph variables are included in the corresponding
blocks. Renaming bound variables and adjoining the algebraic frame
constant uses fewer than twelve blocks of size $O(m+1)$. All atoms
have the format in Lemma~\ref{lem:v63-graph-format}.

Apply Lemma~\ref{lem:v63-effective-modulus}. At a point with $d>0$,
take $h=d$ and $\delta=\min(\chi,d)\ge\chi d$. The definition of
$a_m$ yields \eqref{eq:v63-buffered-separation}; the case $d=0$ is
immediate. On the original source the seam error is at most
$2\varepsilon$. Subtract it from the lower bound for $f$ to obtain
\eqref{eq:v63-rank-threshold}.
\end{proof}

\subsection{A collision-count bound for the tube modulus}
\begin{theorem}[Effective vertical tube length]
\label{thm:v63-effective-tube}
The same budgets, after increasing $C_0$ once, satisfy
\begin{equation}\label{eq:v63-effective-tube}
 \sup_{R\in I}|\{t:\operatorname{dist}((R,t),
                      \widehat{\mathcal C}_{m,\chi/2})<h\}|
       \le K_m\chi^{-1}h^{1/E_m}
       \quad(0<\chi,h\le1).
\end{equation}
An empty reference has zero tube length. In particular both the power
and its constant are now bounded in the collision count and cap.
\end{theorem}
\begin{proof}
Define $b_m(\delta)$ to be the infimum of $\{1\}$ and of the values
$h\in[0,1]$ for which there exist $\delta\le\chi\le1$, $R\in I$
and $u\in[0,mT_0]$ with
\[
 [u,u+\delta]\subset V_{m,\chi/2}(R,h).
\]
It is nondecreasing, since a longer interval contains a shorter one
and the allowed caps decrease when $\delta$ decreases. It is positive.
Indeed, if $h_l\to0$ for a fixed $\delta>0$, pass to limits in
$R_l,\chi_l,u_l$. For each point of $[u_*,u_*+\delta]$, compactness
of the critical graph with cap at least $\delta/2$ places that point
in the limiting fixed-radius critical fiber at cap $\chi_*/2$. This
contradicts its finiteness in Lemma~\ref{lem:v63-graph-format}.

The interval inclusion has the literal formula
\[
 \forall v\in[u,u+\delta]\ \exists Y:\quad
 Y\text{ is a critical graph point at cap }\chi/2,\quad
 |R_Y-R|\le h,\quad F(Y)=v.
\]
Using the lower-bound and approximation descriptions of an infimum
as above gives at most nine quantifier blocks, including the fixed
algebraic coefficient, each with $O(m+1)$ variables. This avoids
forming an integral of a semialgebraic function or quantifying one
variable for every interval endpoint. Lemma~\ref{lem:v63-effective-modulus}
therefore gives $b_m(\delta)\ge K^{-1}\delta^E$ with preliminary
budgets of the form \eqref{eq:v63-budget}.

Let an interval component of $V_{m,\chi/2}(R,h)$ have length $l>0$.
It lies in $[0,mT_0]$. With $\delta=\min(l,\chi,1)$ the preceding
bound gives $h\ge K^{-1}\delta^E$. Since
$\delta\ge\chi l/(1+mT_0)$,
\[
 l\le (1+mT_0)\chi^{-1}(Kh)^{1/E}.
\]
There are at most $CA^m$ interval or point components. The vertical
section of the Euclidean $h$-tube is contained in the $h$-neighborhood
of this set, whose length is at most the sum of these component
lengths plus $2CA^m h$. Absorb the finite factors, $K^{1/E}$ and
$1+mT_0$ by a fixed increase of $C_0$. Increasing $E_m$ weakens the
power bound for $h\le1$, so this change proves
\eqref{eq:v63-effective-tube} simultaneously with the separation
bound.
\end{proof}

The quantitative input in this section is the standard real-algebraic
elimination bound, with its coefficient-height control. Its physical
content is the simple-root graph, the seam/rank finite-fiber argument
and the buffered-cap compactness. The estimates do not independently
certify those inherited continuum inputs. They give explicit, very
large finite-count moduli, not a spectral estimate at a growing band.
